% BevoTag Prior Art Report, memo-style report.
% Formatting follows resources/ECE Style Guide.pdf and resources/Template for Memo Style Reports.doc;
% content follows resources/Prior Art Report description.pdf and rubric (paths relative to this file).
% Preamble matches docs/project-definition/main.tex.
% Build: pdflatex main.tex (twice).
\documentclass[12pt,letterpaper]{article}

\usepackage[T1]{fontenc}
\usepackage{newtxtext,newtxmath}              % 12-pt Times
\usepackage[margin=1in]{geometry}             % 1-in margins on all sides
\usepackage{setspace}
\setstretch{1.43}                             % Word "1.5 lines" for 12-pt Times (20.7-pt baselines)
\usepackage[document]{ragged2e}               % do not justify the right margin
\usepackage{graphicx}
\usepackage{xcolor}
\usepackage{array,tabularx,booktabs}
\usepackage{enumitem}
\usepackage{caption}
\usepackage{float}
\usepackage{flafter}                          % a table never appears before the text that cites it
\usepackage{placeins}
\usepackage{etoolbox}
\usepackage[nobottomtitles*]{titlesec}        % never leave a heading on the last line of a page
\usepackage{varioref}
\usepackage[hidelinks]{hyperref}
\usepackage{url}
\renewcommand{\UrlFont}{\normalfont}
\def\UrlBreaks{\do\/\do\-\do\_\do\=\do\?}  % no break after "." (avoids "doi: 10." line ends)

% Block paragraphs: no indentation, one empty line between paragraphs.
\setlength{\parindent}{0pt}
\setlength{\parskip}{\baselineskip}
\renewcommand{\bottomtitlespace}{0.05\textheight}

% Ragged-right text without hyphenated line breaks.
\hyphenpenalty=10000
\exhyphenpenalty=10000

% Headings (Style Guide 2.3): one empty line before, text starts on the next line, no numbers.
%   1st level: ALL CAPS, bold   2nd level: Initial Caps, bold   3rd level: Initial Caps, bold italic
\setcounter{secnumdepth}{0}
\titleformat{\section}{\normalfont\bfseries}{}{0pt}{\MakeUppercase}
\titleformat{\subsection}{\normalfont\bfseries}{}{0pt}{}
\titleformat{\subsubsection}{\normalfont\bfseries\itshape}{}{0pt}{}
\titlespacing*{\section}{0pt}{0pt}{-\parskip}
\titlespacing*{\subsection}{0pt}{0pt}{-\parskip}
\titlespacing*{\subsubsection}{0pt}{0pt}{-\parskip}

% Captions (Style Guide 2.4): bold, centered, "Table 1." above tables, "Figure 1." below figures.
\captionsetup{font={normalsize,bf},labelfont=bf,labelsep=period,justification=centering,singlelinecheck=false}
\captionsetup[table]{position=above,skip=6pt}
\captionsetup[figure]{position=below,skip=6pt}
\renewcommand{\topfraction}{0.9}
\renewcommand{\bottomfraction}{0.8}
\renewcommand{\textfraction}{0.07}
\renewcommand{\floatpagefraction}{0.9}
\setcounter{topnumber}{3}
\setcounter{totalnumber}{4}
% One empty line between a graphic and the text around it.
\setlength{\textfloatsep}{\baselineskip}
\setlength{\floatsep}{\baselineskip}
\setlength{\intextsep}{\baselineskip}
\makeatletter
\g@addto@macro\float@endH{\vskip-\parskip\vskip\z@}
\makeatother
\setlength{\emergencystretch}{3em}
% No single line of a paragraph stranded at the top or bottom of a page.
\widowpenalty=10000
\clubpenalty=10000

% Indented lists that continue the preceding paragraph (Style Guide 2.5).
\setlist[itemize]{leftmargin=0.5in,topsep=0pt,itemsep=0pt,parsep=0pt}
\setlist[enumerate]{leftmargin=0.5in,topsep=0pt,itemsep=0pt,parsep=0pt}

% Single-spaced 12-pt tables.
\newcolumntype{L}[1]{>{\RaggedRight\arraybackslash}p{#1}}
\newcolumntype{Y}{>{\RaggedRight\arraybackslash}X}
\newenvironment{spectable}{\setstretch{1.0}\renewcommand{\arraystretch}{1.1}}{}

% Page numbers centered at the bottom, on every page but the first.
\pagestyle{plain}

\pretocmd{\section}{\FloatBarrier}{}{}
% Pointer to a graphic that does not share its first reference's page (Style Guide 2.4).
\vrefwarning
\renewcommand{\reftextcurrent}{}
\renewcommand{\reftextafter}{on the next page}
\renewcommand{\reftextfaceafter}{on the next page}
\begin{document}
\thispagestyle{empty}

%----------------------------------------------------------------------------------------------------
% MEMO HEADING (Style Guide 2.2)
%----------------------------------------------------------------------------------------------------
\begingroup
\setstretch{1.0}
\newcommand{\fieldgap}{\\[6.2pt]}
\noindent\begin{tabular}{@{}p{1.1in}@{}p{\dimexpr\linewidth-1.1in\relax}@{}}
\textbf{DATE:} & November 6, 2026 \fieldgap
\textbf{TO:} & Dr.\ Haris Vikalo \\ & Mugdha Jadhao \fieldgap
\textbf{FROM:} & Ethan Chua \\ & Esteban Martinez \\ & Li Shi \\ & Andrew Yin \\ & Abhirit Das \fieldgap
\textbf{SUBJECT:} & Prior Art Report for BevoTag: Wi-Fi Fingerprinting, Inertial Fusion, and
Item-Tracking Patents That Inform the Design \\
\end{tabular}
\endgroup

%----------------------------------------------------------------------------------------------------
\section{Introduction}
%----------------------------------------------------------------------------------------------------
This Prior Art Report documents the BevoTag team's search of published research and patents related
to the project, and explains how the results have shaped the design.
BevoTag is a small tag that locates personal belongings or shared equipment on the University
of Texas at Austin campus by using campus Wi-Fi signals and a motion sensor. The team's Project Definition
committed to this approach but left four questions open: which Wi-Fi measurement to use, how much
surveying a building requires, how to combine Wi-Fi with motion data, and which parts of the design
existing patents already cover. To answer these questions before the system design is finalized,
the team reviewed five research publications and five U.S. patents. This report first summarizes
the design problem and then compares each publication and patent with BevoTag's requirements.
Finally, it explains how the findings have changed the design.

%----------------------------------------------------------------------------------------------------
\section{Overview of the Design Problem}
%----------------------------------------------------------------------------------------------------
Students, faculty, and staff often need to find a misplaced item somewhere on campus, yet no
existing tracker works reliably indoors without help from other people's phones. Satellite
positioning signals, such as those of the Global Positioning System (GPS), do not penetrate
buildings well \cite{hechan}. Consumer tags such as the Apple AirTag and the Samsung Galaxy
SmartTag2 update an item's location only when a nearby phone from the same manufacturer detects
them \cite{airtag}, \cite{smarttag}, so the location of a bag left in an empty lab at night stops
updating. Wi-Fi
positioning that relies on signal strength alone typically errs by 1 to 5~m \cite{rssi}, enough to
report the wrong room. Customers therefore need a tag that reports its own location on any floor,
without a passing phone and without new infrastructure such as Bluetooth beacons installed in
every room.

BevoTag addresses this need by combining two sources of position information. The tag records the
received signal strength indicator (RSSI) of nearby access points and, where its chip supports it,
their channel state information (CSI), which describes how each signal changes across its frequency
subcarriers. The tag also samples an inertial measurement unit (IMU) that measures acceleration and
rotation. A location server compares the Wi-Fi measurements with a survey recorded at known points,
a process called fingerprinting, and fuses the result with the IMU data, because an IMU used alone
produces a position estimate that drifts by more than 150~m within 60~s \cite{imu}. A user
application then shows the position on a campus map.

A successful design will identify the room that contains the item in more than 90\% of trials,
keep the median position error under 5~m, update at least every 10~s while the item moves, and run
at least 24~hours on one charge, at a parts cost of no more than \$50 per tag. The prototype uses
an ESP32-family microcontroller, which exposes CSI to user firmware \cite{espcsi}, and the team has
already built a data logger and a survey application for three buildings, starting with the
Engineering Education and Research Center (EER).

%----------------------------------------------------------------------------------------------------
\section{Review of Prior Art}
%----------------------------------------------------------------------------------------------------
The search was organized around four research questions that the Project Definition left open.
Each question corresponds to one design challenge, and a source was kept only if it gave evidence
on at least one question:
\begin{enumerate}
  \item \textbf{Wi-Fi measurement:} Does CSI give enough accuracy beyond RSSI to justify its cost,
    and can a low-cost chip capture it from campus access points?
  \item \textbf{Survey effort:} How many survey points, samples, and repeated visits does a
    building need, and how quickly does a survey go out of date?
  \item \textbf{Sensor fusion:} How should the tag combine Wi-Fi fixes with IMU data to keep
    drift bounded, and how much does fusion improve accuracy?
  \item \textbf{Freedom to operate and novelty:} Which parts of the design do existing products
    and patents already cover, and where can BevoTag differ?
\end{enumerate}
The first subsection reviews research publications that address questions 1 through 3, and the
second reviews patents that address all four.

\subsection{Prior Art Exclusive of Patent Information}
Five research publications trace Wi-Fi positioning from the first signal-strength fingerprinting
system to CSI capture on the same microcontroller family that BevoTag uses. Two of them establish
RSSI and CSI baselines, one tests CSI capture on BevoTag's hardware, one measures the benefit of
inertial fusion, and one compares many systems in a realistic, changing environment.

\subsubsection{Search Procedure}
The team searched IEEE Xplore with keyword queries that combined ``Wi-Fi'' or ``WLAN'' (wireless
local area network) with ``fingerprint,'' ``channel state information,'' ``ESP32,'' ``pedestrian
dead reckoning,'' ``inertial,'' and ``Kalman'' or ``particle filter,'' and it retrieved each
candidate's full text from the publisher or the authors' own pages. The team kept papers that
report measured accuracy in real buildings, use hardware comparable in cost to BevoTag's, or
evaluate systems in environments that change. Papers were discarded if they relied only on simulation, required
infrastructure that the team cannot install, such as ultra-wideband anchors, or addressed a task
other than locating a device, unless they measured hardware that BevoTag will use.

\subsubsection{RADAR: The First Wi-Fi Fingerprinting System}
Bahl and Padmanabhan's RADAR system, presented in 2000, established the two-phase fingerprinting
method that BevoTag adopts \cite{radar}. In an off-line phase, the authors recorded signal strength
from three base stations at 70 known points on a 43.5~m by 22.5~m office floor with more than 50
rooms; in an on-line phase, they located a laptop by finding the survey points whose signal
strengths were nearest in ``signal space.'' This empirical method achieved a median error of
2.94~m, which the authors noted is about the size of an office room \cite{radar}.

RADAR's results inform BevoTag's survey design more than its accuracy targets. The authors found
that about 40 survey points and three real-time samples achieved nearly all of the accuracy of the
full data set \cite{radar}. When the survey contained only one orientation and the user faced the
opposite way, however, the median error rose 67\% to 4.90~m, because the user's body changes the
signal strength that each base station measures \cite{radar}. RADAR's median error is within
BevoTag's 5~m ceiling but too coarse for reliable room identification, and its three base stations
were dedicated hardware, whereas BevoTag relies on whatever campus access points it hears.

\subsubsection{DeepFi: CSI Fingerprinting Compared with RSSI}
The clearest evidence on whether CSI improves on signal strength comes from Wang, Gao, Mao, and
Pandey's DeepFi system, which compares CSI and RSSI fingerprints measured in the same rooms and was
published in \emph{IEEE Transactions on Vehicular Technology} \cite{deepfi}. The authors recorded
the amplitudes of 30 subcarriers on each of three antennas of a laptop Wi-Fi card, trained a neural
network for each of 50 survey points, and estimated position as a probability-weighted average of
the survey points. They compared DeepFi with an earlier CSI method and with Horus, a widely cited
probabilistic RSSI system \cite{horus}, using the same data.

CSI fingerprints were more accurate than RSSI fingerprints in both test rooms. In a mostly empty
4~m by 7~m living room, DeepFi's mean error was 0.94~m and Horus's was 1.54~m; in a cluttered 6~m
by 9~m laboratory with few line-of-sight paths, the errors were 1.81~m and 2.60~m
\cite{deepfi}. These reductions of 39\% and 30\% justify evaluating CSI for BevoTag. The study
differs from BevoTag in three important ways, however: it used one access point that the authors
controlled, a laptop card with three antennas rather than BevoTag's single-antenna chip, and environments
much smaller than a typical building floor.

\subsubsection{ESP32 CSI Toolkit: CSI Capture on a Microcontroller}
Because DeepFi relied on a laptop card, evidence was also needed that BevoTag's own hardware can
capture CSI. Hernandez and Bulut showed in 2020 that the ESP32 microcontroller can capture CSI
without the laptops and modified drivers that earlier studies required \cite{esp32tool}. Their
open-source toolkit runs on a module that costs less than \$10, weighs less than 10~g, and reports
64 subcarriers from a single antenna \cite{esp32tool}. The authors used it for activity
recognition rather than localization, but they measured the capture rates that matter to BevoTag.

These capture rates reveal a limit on BevoTag's design. When the ESP32 was connected to the
transmitter, it captured CSI from about 600 frames per second; when it only listened, which BevoTag
will do to hear many campus access points, it captured just 34.5\% as many frames, with greater
variation from second to second, because a passive receiver cannot request retransmission
\cite{esp32tool}. This result confirms that the chosen hardware can supply CSI, but it suggests
that the tag will collect far fewer CSI samples from campus access points than from a transmitter
it controls.

\subsubsection{Evennou and Marx: Fusing Wi-Fi with Dead Reckoning}
For the fusion stage, the closest published precedent is a 2006 \emph{EURASIP Journal on Applied
Signal Processing} paper by Evennou and Marx of France Telecom, who fused Wi-Fi fingerprinting with
pedestrian dead reckoning \cite{evennou}. A belt-worn box counted steps with an accelerometer and
tracked heading with a gyroscope. The authors compared a Kalman filter with a particle filter that
uses the building map to eliminate any position hypothesis that crosses a wall.

Each step of filtering and fusion improved accuracy in their 40~m by 40~m office building with
four access points. Along a corridor trajectory, nearest-neighbor fingerprinting gave a mean error
of 3.32~m, a statistical fingerprint method 2.88~m, a Kalman filter 2.56~m, the map-aware particle
filter 1.86~m, and the particle filter with inertial data 1.53~m \cite{evennou}. The authors also
observed that the particle filter's many wall checks make it computationally expensive
\cite{evennou}. These results support fusion and the use of floor plans, which BevoTag's survey
application already contains, but they were measured with a laptop doing the processing and a
survey of only one point per room.

\subsubsection{Microsoft Indoor Localization Competition: Performance Under Change}
Each study above tested one system in a space of the authors' choosing. In contrast, Lymberopoulos
\emph{et al.} reported how 22 indoor positioning systems performed when tested in the same
unfamiliar space at the 2014 Microsoft Indoor Localization Competition \cite{msilc}. Teams were
given seven hours to install hardware or survey an area of about 300~m$^2$ that spanned two rooms
and a hallway, and the organizers then rearranged the furniture in one room before measuring each
system at 20 test points.

The competition shows that Wi-Fi fingerprinting can rival custom infrastructure but degrades when
a room changes. The best system that relied only on existing access points and phone sensors had a
mean error of 1.6~m, better than most systems that installed their own hardware, and half of all
teams achieved errors below 3~m \cite{msilc}. Most teams needed all seven hours to survey or
install, and rearranging the furniture raised the error of the leading Wi-Fi systems by up to
0.94~m \cite{msilc}. Because campus rooms change daily as people move chairs and walk through
hallways, these findings indicate that a one-time survey of EER will not stay accurate.

% Style Guide 2.4: say where a table is when it is not on the citing page; recheck after edits.
Table~\ref{tab:art} compares the five publications with BevoTag. Together, they support
fingerprinting with inertial fusion but leave a single-antenna tag in a changing building untested.

\begin{table}[!htb]
\caption{Research Publications Compared with BevoTag Requirements}
\label{tab:art}
\centering
\begin{spectable}
\begin{tabularx}{\linewidth}{@{}L{1.25in} L{0.75in} Y Y@{}}
\toprule
\textbf{Source} & \textbf{Question} & \textbf{Key Finding} & \textbf{Difference from BevoTag} \\
\midrule
RADAR \cite{radar} & 2 & 2.94~m median; survey should cover different orientations & Dedicated base stations; no motion sensing \\ \addlinespace
DeepFi \cite{deepfi} & 1 & CSI mean error 30--39\% below RSSI & One controlled access point; three-antenna laptop card \\ \addlinespace
ESP32 CSI Toolkit \cite{esp32tool} & 1 & Module under \$10 captures 64-subcarrier CSI; passive capture rate 34.5\% of active & Used for activity sensing, not localization \\ \addlinespace
Evennou and Marx \cite{evennou} & 3 & Map-aware particle filter with inertial data: 1.53~m mean error & Laptop processing; one survey point per room \\ \addlinespace
Microsoft competition \cite{msilc} & 2, 3 & Best infrastructure-free system 1.6~m; furniture changes add up to 0.94~m & Two rooms and a hallway; phones, not tags \\
\bottomrule
\end{tabularx}
\end{spectable}
\end{table}

\subsection{Patent Search and Findings}
Five U.S. patents cover the survey database, the CSI fingerprint, the fusion step, and the
crowdsourced tag model that BevoTag will either build on or avoid. Each patent's claims are
described below and compared with BevoTag's design; these comparisons are engineering judgments,
not legal opinions.

\subsubsection{Search Procedure}
The team searched Google Patents with combinations of ``Wi-Fi'' or ``wireless access point'' with
``fingerprint,'' ``channel state information,'' ``subcarrier,'' ``dead reckoning,'' ``inertial,''
``lost and found,'' and ``tag,'' and it retrieved the full grant documents from the U.S. Patent
and Trademark Office. For each candidate, the team read the independent claims and recorded the
assignee, grant date, and status, keeping the patents whose claims describe a function that
BevoTag performs. Google Patents' status and expiration dates are estimates rather than legal
conclusions, so before the design is finalized, the team will ask the Engineering Library staff to
help confirm the status of any patent of concern.

\subsubsection{US 7,414,988: Server for Updating a Location Beacon Database}
Skyhook Wireless's patent, granted in 2008 to Jones \emph{et al.}, describes a server that
maintains a database of Wi-Fi access point locations across a metropolitan area \cite{skyhook}.
Vehicles carrying a GPS receiver and a Wi-Fi radio drive nearly every street in a target area, and
the server estimates each access point's position from many readings, discards readings that fall
far from a weighted centroid, and updates the database as it discovers new access points. The
patent differs from BevoTag in what it stores: it records where each access point is, whereas
BevoTag records the signal pattern measured at each survey point inside a building. Google Patents
lists the patent's adjusted expiration as September 30, 2026, so its outlier rejection and
database updating are now free to adopt.

\subsubsection{US 9,014,718: Location Sensing Using Channel Fading Fingerprinting}
Intel's patent, granted in 2015 to Yang, Yang, and
Guo, claims the core of a CSI fingerprinting system \cite{intel}. A device scans nearby access
points, measures a ``channel fading profile'' from selected subcarriers, and finds its location by
matching the profile to a database recorded during a survey, using a minimum-distance or
maximum-likelihood rule. The dependent claims allow the database to sit on the device or a server,
which mirrors BevoTag's open question of where to compute the position. This patent is the closest
prior art to BevoTag's Wi-Fi stage, and it means that CSI fingerprint matching by itself is not
patentable. Google Patents lists it as expired because its maintenance fee was not paid, so the
method may be used freely.

\subsubsection{US 9,632,169: Using Inertial Navigation to Calibrate an Indoor Positioning System}
Ricoh's patent, granted in 2017 to Moraleda and Hull, addresses survey effort \cite{ricoh}. A
surveyor walks between known points while inertial navigation tracks the path and Wi-Fi
signatures are collected continuously; each signature is then assigned a position along a B-spline
fitted to the path, and the surveyor may adjust the spline's control points on a display. This
method could survey a floor far faster than stopping at each point. The patent remains in force
until 2035, so the team will avoid the method of fitting and editing a spline path.

\subsubsection{US 9,332,476: Correcting Indoor Positioning Using Wi-Fi Handovers}
BlackBerry's patent, granted in 2016 to Shatsky and Kouznetsov, uses Wi-Fi to correct dead-reckoning 
drift when a device hands over to a new access point \cite{blackberry}. If the dead-reckoning
position lies outside a ``handover spot'' around the new access point, the method of claim~1 resets the 
position to the nearest point on that spot's border. This patent confirms that drift correction from Wi-Fi
is a recognized need, but it corrects only at handovers and only coarsely. BevoTag will instead correct
the inertial estimate at every fingerprint fix with a filter, the approach that Evennou and Marx published in
2006 \cite{evennou}, and will not reset the estimate to a handover boundary. The patent remains in
force until 2034.

\subsubsection{US 9,479,920: Power Management in a Crowdsourced Lost-and-Found Service}
Finally, Apple's patent, granted in 2016 to Mayor, describes the crowdsourced model behind tags
such as the AirTag \cite{apple}. A Bluetooth Low Energy tag broadcasts an identifier; low-power
subsystems in a passing phone record that identifier and the Wi-Fi access points the phone hears,
each with a time stamp; and the phone's main processor periodically wakes, assigns the tag the
stored location of the access point heard at the same time, and reports it to a server. The patent
locates a tag only as precisely as the nearest access point and only when a participating phone is
nearby, which are the two weaknesses that BevoTag is designed to overcome. Because BevoTag scans
and reports by itself and uses no helper phones, it does not practice these claims, but the
patent's strategy of batching measurements to keep the main processor asleep applies directly to
BevoTag's battery target.

% Style Guide 2.4: say where a table is when it is not on the citing page; recheck after edits.
Table~\ref{tab:pat} summarizes the five patents and their status: the two expired patents are free
to use, and BevoTag's design avoids the claims of the three active ones.

\begin{table}[!htb]
\caption{Patents Relevant to BevoTag}
\label{tab:pat}
\centering
\begin{spectable}
\begin{tabularx}{\linewidth}{@{}L{1.05in} L{1.0in} Y L{1.2in}@{}}
\toprule
\textbf{Patent} & \textbf{Assignee} & \textbf{Relevance to BevoTag} & \textbf{Status (Google Patents)} \\
\midrule
US 7,414,988 \cite{skyhook} & Skyhook Wireless & Access point database with outlier rejection and updating & Expired Sept.~2026 \\ \addlinespace
US 9,014,718 \cite{intel} & Intel & CSI fingerprint matching against a survey database & Expired (unpaid fee) \\ \addlinespace
US 9,632,169 \cite{ricoh} & Ricoh & Inertially tracked walking survey fitted with a B-spline & Active to 2035; avoid \\ \addlinespace
US 9,332,476 \cite{blackberry} & BlackBerry & Drift reset to a handover-spot border & Active to 2034; avoid \\ \addlinespace
US 9,479,920 \cite{apple} & Apple & Crowdsourced Bluetooth tag located by phones & Active to 2035; not practiced \\
\bottomrule
\end{tabularx}
\end{spectable}
\end{table}

%----------------------------------------------------------------------------------------------------
\section{Impact of Prior-Art Search on Design Decision-Making}
%----------------------------------------------------------------------------------------------------
The search answered each of the four research questions and changed three parts of the design
plan: the treatment of CSI, the survey protocol, and the fusion design. The following subsections
answer each question in turn and state the decision that follows.

\subsection{Wi-Fi Measurement: CSI Is Promising but Unproven on BevoTag's Hardware}
The literature supports CSI's accuracy but leaves its feasibility on BevoTag in doubt: DeepFi's 30
to 39\% error reduction \cite{deepfi} came from hardware far more capable than BevoTag's
single-antenna tag, which will capture far fewer CSI frames whenever it listens passively
\cite{esp32tool}.

The team therefore decided to make RSSI fingerprinting the baseline for every later design and to
evaluate CSI as an addition to it. The first data logger experiment will record CSI and RSSI
together in EER and measure how many usable CSI frames per second the tag captures from campus
access points, both while associated with the campus utexas-iot network and while listening
passively. If CSI improves room identification on held-out survey points by a margin comparable to
DeepFi's, the design will keep it; otherwise, the design will drop it to save processing and power.
Because Intel's CSI fingerprinting patent has expired \cite{intel}, either outcome is free of
patent risk.

\subsection{Survey Effort: Orientation and Time Matter More than Density}
RADAR and the competition together show that a survey's coverage of conditions matters more than
its density. RADAR reached nearly full accuracy with about 40 points, but its median error rose
67\% when the survey and the user faced opposite directions \cite{radar}, and the competition's
leading Wi-Fi systems gained up to 0.94~m of error when furniture moved \cite{msilc}.

The team has accordingly revised the survey protocol in four ways:
\begin{enumerate}
  \item Scans will be recorded facing at least four directions at each survey point.
  \item Each benchmark area will be surveyed at three times of day, including class changes, when
    hallways are crowded.
  \item The survey points from one visit will be held out for testing, so that accuracy figures
    reflect a survey that has aged.
  \item The database will be kept current without full resurveys by adapting the outlier rejection
    and incremental updating of Skyhook's now-expired patent \cite{skyhook}.
\end{enumerate}

\subsection{Sensor Fusion: A Map-Aware Filter Bounds Drift}
Evennou and Marx's results show that fusion and floor-plan constraints together roughly halve
fingerprinting error, from 3.32~m to 1.53~m \cite{evennou}, which suggests that BevoTag's 5~m
median target is conservative and that room identification is the harder goal. Their study, the
competition's mix of Wi-Fi and inertial entries \cite{msilc}, and BlackBerry's patent
\cite{blackberry} all treat Wi-Fi as the source that corrects inertial drift, confirming the
architecture in the Project Definition.

The fusion design has therefore narrowed to two candidate filters. An extended Kalman filter
\cite{welch} is light enough to run on the tag but cannot easily use walls, whereas a particle
filter can discard paths through walls using the floor plans already in the survey application but
is computationally expensive \cite{evennou}. Because of this trade-off, a particle filter would
most likely run on the server, and the tag would then send measurements rather than positions. The
team will compare the two filters on logged EER walks before deciding.

\subsection{Freedom to Operate and Novelty: The Combination Is the Contribution}
The search shows that each of BevoTag's building blocks is already public: the patents in
Table~\ref{tab:pat} claim most of them, and the research literature has published Wi-Fi and inertial
fusion since at least 2006 \cite{evennou}. No source found in the search, however, combines
CSI-capable fingerprinting and inertial fusion on a self-contained, battery-powered tag that reports
its own position without helper phones or new infrastructure. BevoTag will therefore build on the
two expired patents and published methods, avoid the specific claims of the three active patents,
and rest any claim of novelty on the integrated tag and its performance on a changing campus rather
than on a new positioning algorithm.

%----------------------------------------------------------------------------------------------------
\section{Conclusion}
%----------------------------------------------------------------------------------------------------
This report presented a prior-art search that answered four open questions about BevoTag's Wi-Fi
measurement, survey, sensor fusion, and patent landscape. The publications show that CSI can reduce
error but that the tag will capture it less often, that a survey needs to cover several
orientations and times of day, and that a map-aware fusion filter can roughly halve position error;
the patents show that two have expired and that the design avoids the claims of the other three. RSSI fingerprinting will therefore serve as the baseline, CSI as an experiment, and the
integrated, self-reporting tag as the project's contribution. Next, the team will measure CSI
capture in EER, survey a benchmark area under the revised protocol, and compare two fusion filters
to decide where BevoTag computes its position.

%----------------------------------------------------------------------------------------------------
% REFERENCES (IEEE style, Style Guide 3.2): single-spaced entries, an empty line between entries,
% numbered in order of first citation, bracketed numbers flush left with the text at a 0.5-in tab.
%----------------------------------------------------------------------------------------------------
\renewcommand{\refname}{References}
\makeatletter
\renewenvironment{thebibliography}[1]
  {\section{\refname}%
   \setstretch{1.0}%
   \list{\@biblabel{\@arabic\c@enumiv}}%
        {\settowidth\labelwidth{\@biblabel{#1}}%
         \setlength{\leftmargin}{0.5in}%
         \setlength{\labelsep}{\dimexpr0.5in-\labelwidth\relax}%
         \setlength{\itemsep}{\baselineskip}%
         \setlength{\parsep}{0pt}%
         \setlength{\topsep}{0pt}%
         \usecounter{enumiv}%
         \renewcommand\theenumiv{\@arabic\c@enumiv}%
         \def\makelabel##1{##1\hfil}}%
   \RaggedRight\exhyphenpenalty=50 }
  {\endlist}
\makeatother

\begin{thebibliography}{10}
\bibitem{hechan}
S. He and S.-H. G. Chan, ``Wi-Fi fingerprint-based indoor positioning: Recent advances and
comparisons,'' \emph{IEEE Commun. Surveys Tuts.}, vol.~18, no.~1, pp.~466--490, 1st Quart. 2016,
doi: \url{10.1109/COMST.2015.2464084}.

\bibitem{airtag}
Apple Inc. (2026). \emph{AirTag} [Online]. Available: \url{https://www.apple.com/airtag/}.

\bibitem{smarttag}
Samsung Electronics Canada. (n.d.). \emph{Galaxy SmartTag2} [Online]. Available:
\url{https://www.samsung.com/ca/business/mobile-accessories/galaxy-smarttag2-white-ei-t5600bwegca/}.

\bibitem{rssi}
A. S. Yaro \emph{et al.}, ``A survey of the performance-limiting factors of a
2-dimensional RSS fingerprinting-based indoor wireless localization system,'' \emph{Sensors},
vol.~23, no.~5, art.~no.~2545, Feb. 2023, doi: \url{10.3390/s23052545}.

\bibitem{imu}
O. J. Woodman, ``An introduction to inertial navigation,'' Comput. Lab., Univ. Cambridge,
Cambridge, U.K., Tech. Rep. UCAM-CL-TR-696, Aug. 2007, doi: \url{10.48456/tr-696}.

\bibitem{espcsi}
Espressif Systems. (2026). \emph{ESP-IDF Programming Guide (ESP32-S3): Wi-Fi Vendor Features}
[Online]. Available:
\url{https://docs.espressif.com/projects/esp-idf/en/stable/esp32s3/api-guides/wifi-driver/wifi-vendor-features.html}.

\bibitem{radar}
P. Bahl and V. N. Padmanabhan, ``RADAR: An in-building RF-based user location and tracking
system,'' in \emph{Proc. IEEE INFOCOM 2000}, Tel Aviv, Israel, Mar. 2000, vol.~2, pp.~775--784,
doi: \url{10.1109/INFCOM.2000.832252}.

\bibitem{deepfi}
X. Wang \emph{et al.}, ``CSI-based fingerprinting for indoor localization: A deep learning
approach,'' \emph{IEEE Trans. Veh. Technol.}, vol.~66, pp.~763--776, Jan. 2017,
doi: \url{10.1109/TVT.2016.2545523}.

\bibitem{horus}
M. Youssef and A. Agrawala, ``The Horus WLAN location determination system,'' in \emph{Proc. 3rd
Int. Conf. Mobile Syst., Appl., Services (MobiSys)}, Seattle, WA, USA, Jun. 2005, pp.~205--218,
doi: \url{10.1145/1067170.1067193}.

\bibitem{esp32tool}
S. M. Hernandez and E. Bulut, ``Lightweight and standalone IoT based WiFi sensing for active
repositioning and mobility,'' in \emph{Proc. IEEE 21st Int. Symp. World Wireless, Mobile Multimedia
Netw. (WoWMoM)}, Cork, Ireland, Aug.--Sept. 2020, pp.~277--286,
doi: \url{10.1109/WoWMoM49955.2020.00056}.

\bibitem{evennou}
F. Evennou and F. Marx, ``Advanced integration of WiFi and inertial navigation systems for indoor
mobile positioning,'' \emph{EURASIP J. Appl. Signal Process.}, vol.~2006, art.~no.~86706,
pp.~1--11, 2006, doi: \url{10.1155/ASP/2006/86706}.

\bibitem{msilc}
D. Lymberopoulos \emph{et al.}, ``Microsoft indoor localization competition: Experiences and
lessons learned,'' \emph{GetMobile: Mobile Comput.
Commun.}, vol.~18, no.~4, pp.~24--31, Oct. 2014, doi: \url{10.1145/2721914.2721923}.

\bibitem{skyhook}
R. K. Jones \emph{et al.}, ``Server for updating location beacon database,'' U.S. Patent 7\,414\,988, Aug. 19, 2008.

\bibitem{intel}
L. T. Yang \emph{et al.}, ``Location sensing using channel fading fingerprinting,'' U.S.
Patent 9\,014\,718, Apr. 21, 2015.

\bibitem{ricoh}
J. Moraleda and J. J. Hull, ``Using inertial navigation for calibration of indoor positioning
system,'' U.S. Patent 9\,632\,169, Apr. 25, 2017.

\bibitem{blackberry}
A. Shatsky and A. Kouznetsov, ``Method and apparatus to correct indoor positioning by utilizing
Wi-Fi handovers,'' U.S. Patent 9\,332\,476, May 3, 2016.

\bibitem{apple}
R. Mayor, ``Power management in crowd-sourced lost-and-found service,'' U.S. Patent
9\,479\,920, Oct. 25, 2016.

\bibitem{welch}
G. Welch and G. Bishop, ``An introduction to the Kalman filter,'' Dept. Comput. Sci., Univ. North
Carolina, Chapel Hill, Tech. Rep. TR 95-041, Apr. 2004.

\end{thebibliography}

\end{document}
